\chapter{Use case}
\label{ch:use-case}

\section{Sáu use case}

\begin{table}[htbp]
\centering
\caption{Sáu use case và actor chính của từng cái}
\label{tab:sau-use-case}
\begin{tabular}{@{}llL{5.4cm}@{}}
\toprule
\textbf{Mã} & \textbf{Actor chính} & \textbf{Tên} \\
\midrule
UC-01 & Teacher & Tạo đề nháp \\
UC-02 & Teacher & Review và phát hành đề \\
UC-03 & Student & Làm và nộp bài (giai đoạn 1) \\
UC-04 & Student & Chữa bài sau khi nộp (giai đoạn 2) \\
UC-05 & --- & Hệ thống tạo câu biến thể \\
UC-06 & Student & Báo cáo chỗ trợ lí giải thích chưa rõ \\
\bottomrule
\end{tabular}
\end{table}

UC-05 không có actor trực tiếp: đó là hành vi hệ thống, được \texttt{<<include>>} từ UC-04.

\begin{landscape}
\begin{figure}[p]
  \centering
  \includegraphics[width=\linewidth,height=\textheight,keepaspectratio]{use-case.pdf}
  \caption[Use Case Diagram]{Use Case Diagram}
  \label{fig:use-case}
\end{figure}
\end{landscape}

\section{Đặc tả từng use case}

\subsection{UC-01 --- Tạo đề nháp}

\paragraph{Điều kiện trước.} Teacher biết chủ đề hoặc mục tiêu học tập cần kiểm tra, và mô tả
được yêu cầu ở mức đủ hiểu.

\paragraph{Luồng chính.} Teacher nhập yêu cầu và các ràng buộc (chủ đề, số câu, độ khó, thời gian
làm bài, loại bài). Hệ thống phân tích, sinh đề nháp gồm câu hỏi, phương án, đáp án đúng, lời giải
và metadata, rồi hiển thị cho Teacher.

\paragraph{Quan hệ.} \emph{Tạo đề nháp} \texttt{<<include>>} \emph{Sinh nội dung đề}. \emph{Hỏi bổ
sung yêu cầu} \texttt{<<extend>>} \emph{Tạo đề nháp}, khi yêu cầu ban đầu chưa đủ rõ --- đây chính
là một trong hai nơi Giáo viên ra quyết định ở mục~\ref{sec:hai-cong}.

\subsection{UC-02 --- Review và phát hành đề}

\paragraph{Điều kiện trước.} Đã có đề nháp.

\paragraph{Luồng chính.} Teacher xem từng câu hỏi, xem đáp án đúng, lời giải và phương án sai;
chỉnh sửa nếu cần; duyệt đề; rồi phát hành với sáu tham số. Mỗi lần phát hành cho một lớp
sinh ra một \emph{Publication}.

\paragraph{Quan hệ.} \emph{Tạo lại nội dung đề} \texttt{<<extend>>} \emph{Review và phát hành đề},
khi Teacher thấy một phần hoặc toàn bộ đề chưa phù hợp.

\subsection{UC-03 --- Làm và nộp bài}

\paragraph{Điều kiện trước.} Bài đã được phát hành và Student có quyền truy cập.

\paragraph{Luồng chính.} Student mở bài, đọc câu hỏi, chọn đáp án, nộp bài. Hệ thống ghi nhận, và chấm bài.

\paragraph{Quan hệ.} \emph{Làm bài thường xuyên} và \emph{Làm bài cuối kỳ} đều generalize lên
\emph{Làm và nộp bài (giai đoạn 1)}.

Sơ đồ \emph{không} dùng \texttt{<<include>>} từ \emph{Làm và nộp bài} sang \emph{Chấm bài}: 
chấm bài là bước xử lý \emph{sau} khi có bài nộp, không phải hành vi bắt buộc để Student hoàn thành mục tiêu nộp bài của mình.

\paragraph{Kết quả.} Bài nộp được ghi nhận và được chấm. \textbf{Nộp bài không kết thúc bài kiểm
tra} --- nó kết thúc giai đoạn 1.

\subsection{UC-04 --- Chữa bài}

\paragraph{Điều kiện trước.} Student đã nộp bài ở giai đoạn 1 và hệ thống đã chấm; chưa quá hạn kết
thúc giai đoạn 2.

\paragraph{Luồng chính.} Student mở kết quả; mở màn hỏi trợ lí, nơi mọi câu sai nằm cùng một màn
và phần giải thích không tính giờ; bấm nút làm bài mới để khởi động đồng hồ của lượt; làm câu biến
thể của từng câu sai. Câu đúng thì chốt 0,5; câu còn sai thì sang vòng tiếp, tối đa ba
vòng, hết ba vòng thì chốt 0.

\paragraph{Ngoại lệ.} Hết hạn giai đoạn 2 giữa một lượt thì lượt bị \emph{dừng} và các câu còn dở nhận
0 điểm. Student không mở giai đoạn 2 lần nào thì mọi câu sai nhận 0 điểm.

\paragraph{Quan hệ.} \emph{Chữa bài} \texttt{<<include>>} \emph{Chấm bài và tra cứu lỗi sai}, và
\texttt{<<include>>} \emph{Làm câu biến thể, tối đa ba vòng}. \emph{Yêu cầu giải thích cách làm}
và \emph{Báo cáo giải thích chưa rõ} đều \texttt{<<extend>>} từ nó.

Quan hệ với \emph{Làm câu biến thể} là \texttt{<<include>>} chứ không phải \texttt{<<extend>>},
\textbf{vì giai đoạn 2 bắt buộc}: nó không phải một nhánh có thể không xảy ra. Đây là ví dụ rõ nhất
trong cả sơ đồ về việc một quyết định sư phạm đổi luôn kiểu quan hệ trên hình.

\subsection{UC-05 --- Tạo câu biến thể}

\paragraph{Điều kiện trước.} Một câu ở giai đoạn 1 bị làm sai.

\paragraph{Luồng chính.} Hệ thống lấy câu Student làm sai, sinh biến thể của \emph{chính câu
đó} --- giữ nguyên cấu trúc và lỗi cần kiểm, thay dữ kiện --- rồi đưa thẳng tới Student. Kết quả của
câu biến thể quyết định câu gốc chốt 0,5 hay sang vòng tiếp.

\subsection{UC-06 --- Báo cáo Trợ lí giải thích chưa rõ}

\paragraph{Điều kiện trước.} Student đang ở giai đoạn 2, và trợ lí đã trả lời ít nhất một lần.

\paragraph{Luồng chính.} Student bấm nút báo cáo ``Trợ lí giải thích chưa rõ''. 
Hệ thống ghi lại \textbf{cả đoạn hội thoại của bài kiểm tra đó} kèm những câu Student
làm sai để giáo viên có thể xem được.
